\clearpage\begingroup\SourceLeaf{398b}\ArithmeticHead{67}{R}{DIE KOMPLEXEN GRÖSSEN IN IHRER BEZIEHUNG ZUR GEOMETRIE.}
\renewcommand{\thefootnote}{\arabic{footnote})}\setcounter{footnote}{0}
\begin{center}
\textbf{20.}\par\smallskip
\textls[100]{Komplexe Grössen mit mehr als zwei Einheiten.}
\end{center}
\noindent\hspace*{\parindent}In dem Brief an \textsc{Grassmann} vom 14. Dezember 1844 sagt \textsc{Gauss}, auf
dessen ihm übersandte Ausdehnungslehre Bezug nehmend, »dass die Tendenzen
derselben \textls[100]{teilweise} denjenigen Wegen begegnen, auf denen ich selbst nun
fast seit einem halben Jahrhundert gewandelt bin und wovon freilich nur ein
kleiner Teil 1831 in den Comment. der Göttingischen Societät und noch mehr
in den Göttingischen Gelehrten Anzeigen (1831, Stück 64) gleichsam im Vorbeigehen
erwähnt ist; nämlich die konzentrierte Metaphysik der komplexen Grössen,
während von der unendlichen Fruchtbarkeit dieses Prinzips für Untersuchungen
räumliche Verhältnisse betreffend zwar vielfältig in meinen Vorlesungen gehandelt[\footnote{Zum Beispiel hat \textsc{Gauss} vom Dezember 1839 bis Ostern 1740 eine Vorlesung \textit{über die Theorie der imaginären Grössen} gehalten, von der zwei Stücke in den Werken abgedruckt sind (W. VIII, S. 331—334 und S. 346—347).}],
aber Proben davon nur hin und wieder, und \textls[100]{als solche} nur dem
aufmerksamen Auge erkennbar, bei andern Veranlassungen mitgeteilt sind«
(W. X\textsubscript{1}, S. 436). Solche Proben finden sich in der Dissertation, in der Kopenhagener
Preisschrift und in verschiedenen kleineren Aufsätzen zur elementaren
Mathematik, über die im vierten Abschnitt berichtet werden wird.
\par
Von der Selbstanzeige in den Göttingischen Anzeigen kommt hier besonders
der Schluss in Betracht. »Der Verfasser hat sich vorbehalten, den
Gegenstand [der komplexen Grössen], welcher in der vorliegenden Abhandlung
\par\smallskip\noindent\hfill{\small 9\textsuperscript{*}}\par
\clearpage\SourceLeaf{399}\ArithmeticHead{68}{L}{STÄCKEL, GAUSS ALS GEOMETER.}\setcounter{footnote}{0}
\noindent eigentlich nur gelegentlich berührt ist, künftig vollständiger zu bearbeiten, wo
dann auch die Frage, warum die Relationen zwischen Dingen, die eine Mannigfaltigkeit
von mehr als zwei Dimensionen darbieten, nicht noch andere in
der allgemeinen Arithmetik zulässige Grössen liefern können, ihre Beantwortung
finden wird« (W. II, S. 178).
\par
Leider ist \textsc{Gauss} nicht dazu gekommen, das hier gegebene Versprechen
einzulösen, und auch die wenigen im Nachlass vorhandenen Aufzeichnungen,
die man damit in Beziehung bringen kann, reichen nicht aus, um festzustellen,
was er mit seinen Andeutungen gemeint hatte.
\par
Ebenso wie den Punkten der Ebene aus den Einheiten \(1\) und \(i\) gebildete
\textls[100]{bikomplexe} Grössen (W. VIII, S. 354) zugeordnet werden, kann man für
die Punkte des Raumes \textls[100]{trikomplexe} Grössen benutzen (W. VIII, S. 353,
354). Gelegentlich hat \textsc{Gauss} geradezu den drei kartesischen Koordinaten
\(x,y,z\) die drei Einheiten \(1,i,k\) zugesellt und zum Beispiel die Ecken eines
Ikosaeders und eines Dodekaeders durch trikomplexe Grössen \(x+iy+kz\) dargestellt
(Handbuch 16 Bb, S. 166). Es entsteht dann die Frage, wie man mit
solchen Grössen rechnen und im besonderen, wie man für sie das Produkt definieren
soll. \textsc{Gauss} hat, den Kern des Problems erfassend, schon 1819 viergliedrige
komplexe Grössen betrachtet, die er Mutationsskalen nennt (W. VIII,
S. 357—362). Ihre geometrische Bedeutung besteht darin, dass sie die Drehung
eines Raumes in einem andern Raume verbunden mit einer Vergrösserung
oder Verkleinerung ausdrücken, und \textsc{Gauss} ist dazu gelangt, die Multiplikation
zweier solcher Grössen so zu erklären, dass das Produkt das geometrische
Ergebnis zweier hintereinander ausgeführter Mutationen darstellt. Auf diese
Art ist er zu einem Multiplikationsgesetz gelangt, das mit dem \textsc{Hamilton}schen
Quaternionen übereinstimmt.
\par
Weitere Ausführungen über die mehrdimensionalen Mannigfaltigkeiten
bei \textsc{Gauss} findet man in Nr. 33 dieses Aufsatzes.
\clearpage\endgroup
